\section{Experimental Setup, Baselines, and Stronger Models}

\subsection{Evaluation stages}

We evaluate in cumulative stages rather than through a single leaderboard:
\begin{itemize}[leftmargin=1.5em]
\item judge-style pilots for diagnosis,
\item same-backbone frozen-head baselines,
\item minimal repair variants,
\item stronger rerankers,
\item robustness and reproduction layers.
\end{itemize}
This ordering reflects the central question. We first establish
whether answer-visible verification is shortcut-heavy at all; only then do we
ask whether lightweight repairs improve that behavior.

\subsection{Judge-style evaluators}

We first test local prompt judges and an external API judge. These are useful
because they resemble how process evaluation is often deployed in practice, but
they are weak scientific instruments for local process grounding. Across easy,
hard, and naturalized slices, they show the same qualitative pattern: high
answer sensitivity and weak \amcd. This establishes a distinction that
persists throughout: judge-style leakage is real, but it should not
be conflated with the behavior of trained verifiers.

\subsection{Same-backbone frozen-head baselines}

The next stage uses a frozen \texttt{Qwen3-1.7B} backbone with small learned
heads:
\begin{itemize}[leftmargin=1.5em]
\item \texttt{visible\_bce},
\item \texttt{masked\_bce},
\item \texttt{pairwise\_visible},
\item \texttt{step\_only}.
\end{itemize}
These baselines are the first place where the strong masked baseline becomes
decisive. On early synthetic slices, the answer-masked baseline is already
very strong. This result prevents us from overstating the necessity of more
complex repair objectives.

\subsection{Minimal repairs}

We test minimal repair variants:
\begin{itemize}[leftmargin=1.5em]
\item \texttt{local-pair},
\item \texttt{cond-swap},
\item \texttt{hard-neg},
\item minimal and step-scan dual-head variants.
\end{itemize}
\texttt{visible\_cond\_swap} can improve faithfulness and sometimes \amcd,
especially on structured hard OOD slices. However, no simple repair variant
uniformly dominates the masked baseline across naturalized OOD, and the
dual-head line remains negative at the current scale. The resulting picture is
a faithfulness-utility frontier rather than a universal repair win.

\subsection{Stronger rerankers}

We then move beyond frozen linear heads to stronger local rerankers:
\begin{itemize}[leftmargin=1.5em]
\item \texttt{Qwen3-Reranker-8B},
\item \texttt{Qwen3-Reranker-4B}.
\end{itemize}
These models change the deployment picture materially. The strongest current
model is \texttt{Qwen3-Reranker-8B} in the masked setting, not a visible
reranker and not a frozen-head repair.

\subsection{Reporting regimes}

Unless otherwise stated, the main table uses \textsc{Craft-Deploy} because it
best matches the deployment question. \textsc{Craft-Clean} is always
reported as a secondary cleaner benchmark, not as a silent replacement for the
main deployment regime. The two benchmarks separate process-faithfulness and
deployment utility more clearly than either one does alone. We also report an
external-source corroboration line on \textsc{ProcessBench}. This line is not
used for quartet metrics; instead, it checks whether the broader audit story
survives on real step-level supervision without self-generated source
problems. In the released package, this external-source line is pinned through
a canonical suite summary, manifest, and split manifest, so it should be read
as a formal corroboration benchmark rather than as a loose collection of
auxiliary experiments.
